\subsection{Complete numerical realization statement}
The following collects the intrinsic and effective conclusions. Its construction and resource bounds are proved in the surrounding appendices; no stronger numerical converse is implicit.
\begin{theorem}[Sharp causal memory and certified numerical realization]
\label{thm:effective-main}
Under the preceding fixed-experiment hypotheses, with $r\ge2$, the
checkpoint and causal regret laws have orders $\Psi_{n,m}$ and $\Xi_N$,
respectively, uniformly in $a\in K$ and integers $M\ge1$.
Suppose the command cube is rationally specified and certified finite
approximations to the cell coefficients and prior moments through formal
degree $N$ are available at requested tolerances, together with coefficient and
positivity bounds.
An algorithm constructs integer transition tables and dyadic output
tables with at most $M$ persistent labels at each stage and maximum
worst-history regret at most $C\Xi_N(M,a)$. For every such filter with
a fixed memoryless command-coding rule, the maximum unconditional
checkpoint regret under the common exploration law is at least
$c\Xi_N(M,a)$.

The algorithm requires no collision stratum, additive-gap lower bound,
or value of $\Xi_N$ as input. Its selected numerical precision satisfies
\[
 b(M,a)=\tfrac12\log_2(1/\Xi_N(M,a))+O(1),
\]
and a sufficient read-only program size is
$O(M\Xi_N(M,a)^{-J/2}\log(M+1))$ bits. In particular, setting
$d_0=(r-1)\lfloor N/2\rfloor$, one may uniformly use
\[
 b(M)\le d_0^{-1}\log_2(M+1)+O(1),\qquad
 P(M)=O(M^{1+J/d_0}\log(M+1)).
\]
These numerical bounds are sufficient; no precision or program-length
converse is asserted. The running filter retains only its index, not a
posterior vector, historical tape, or numerical oracle. For computable
numerical inputs the synthesis is uniform and terminating. All constants
may depend on the fixed prior and horizon, but not on $a$ or $M$.
\end{theorem}
